
\documentclass{article}
\usepackage{colortbl}
\usepackage{makecell}
\usepackage{multirow}
\usepackage{supertabular}

\begin{document}

\newcounter{utterance}

\centering \large Interaction Transcript for game `cladder', experiment `full\_v1.5\_default', episode 2287 with reflexion{-}Qwen3.5{-}4B.
\vspace{24pt}

{ \footnotesize  \setcounter{utterance}{1}
\setlength{\tabcolsep}{0pt}
\begin{supertabular}{c@{$\;$}|p{.15\linewidth}@{}p{.15\linewidth}p{.15\linewidth}p{.15\linewidth}p{.15\linewidth}p{.15\linewidth}}
    \# & \multicolumn{2}{c}{Player} && \multicolumn{2}{c}{Game Master} \\
    \hline

    \theutterance \stepcounter{utterance}  
    & & & \multicolumn{4}{p{0.6\linewidth}}{
        \cellcolor[rgb]{0.9,0.9,0.9}{
            \makecell[{{p{\linewidth}}}]{
                \texttt{\tiny{[P1$\langle$GM]}}
                \texttt{You are an expert in causal inference. The following question is not a typical commonsense query, but rather a meticulously designed question created by a professor specializing in causal inference, intended to assess the students' mastery of the course content.} \\
\texttt{Imagine a self{-}contained, hypothetical world with only the following conditions, and without any unmentioned factors or causal relationships: Smoking has a direct effect on effort and college admission. Effort has a direct effect on college admission. The overall probability of smoking is 93\%. The probability of nonsmoking and college admission is 3\%. The probability of smoking and college admission is 62\%. Is the chance of college admission smaller when observing smoking?} \\
\texttt{Start your answer with 'Yes' or 'No', followed by additional reasoning or evidence to support your explanation.} \\
            }
        }
    }
     \\ \\

    \theutterance \stepcounter{utterance}  
    & \multicolumn{4}{p{0.6\linewidth}}{
        \cellcolor[rgb]{0.561,0.737,0.561}{
            \makecell[{{p{\linewidth}}}]{
                \texttt{\tiny{[P1$\rangle$GM]}}
                \texttt{Yes} \\
\\ 
\texttt{**Reasoning:**} \\
\\ 
\texttt{To determine whether the chance of college admission is smaller when observing smoking, we must compare the conditional probability of admission given smoking (\$P(\textbackslash{}text\{Admission\} | \textbackslash{}text\{Smoking\})\$) with the conditional probability of admission given nonsmoking (\$P(\textbackslash{}text\{Admission\} | \textbackslash{}text\{Nonsmoking\})\$).} \\
\\ 
\texttt{We are given the following joint probabilities:} \\
\texttt{1.  \$P(\textbackslash{}text\{Smoking\}) = 0.93\$} \\
\texttt{2.  \$P(\textbackslash{}text\{Nonsmoking\}) = 1 {-} 0.93 = 0.07\$} \\
\texttt{3.  \$P(\textbackslash{}text\{Smoking\} \textbackslash{}cap \textbackslash{}text\{Admission\}) = 0.62\$} \\
\texttt{4.  \$P(\textbackslash{}text\{Nonsmoking\} \textbackslash{}cap \textbackslash{}text\{Admission\}) = 0.03\$} \\
\\ 
\texttt{First, we calculate the probability of admission given smoking:} \\
\texttt{\$\$P(\textbackslash{}text\{Admission\} | \textbackslash{}text\{Smoking\}) = \textbackslash{}frac\{P(\textbackslash{}text\{Smoking\} \textbackslash{}cap \textbackslash{}text\{Admission\})\}\{P(\textbackslash{}text\{Smoking\})\} = \textbackslash{}frac\{0.62\}\{0.93\} \textbackslash{}approx 0.6667 \textbackslash{}text\{ (or \} 66.67\textbackslash{}\%)\$\$} \\
\\ 
\texttt{Next, we calculate the probability of admission given nonsmoking:} \\
\texttt{\$\$P(\textbackslash{}text\{Admission\} | \textbackslash{}text\{Nonsmoking\}) = \textbackslash{}frac\{P(\textbackslash{}text\{Nonsmoking\} \textbackslash{}cap \textbackslash{}text\{Admission\})\}\{P(\textbackslash{}text\{Nonsmoking\})\} = \textbackslash{}frac\{0.03\}\{0.07\} \textbackslash{}approx 0.4286 \textbackslash{}text\{ (or \} 42.86\textbackslash{}\%)\$\$} \\
\\ 
\texttt{**Comparison:**} \\
\texttt{The probability of admission for smokers (\$\textbackslash{}approx 66.7\textbackslash{}\%\$) is strictly greater than the probability of admission for nonsmokers (\$\textbackslash{}approx 42.9\textbackslash{}\%\$).} \\
\\ 
\texttt{Therefore, observing smoking is associated with a **higher** chance of college admission in this specific hypothetical world, not a smaller one. The answer to the question "Is the chance of college admission smaller when observing smoking?" is **No**.} \\
\\ 
\texttt{*(Note: While the prompt asks to start with "Yes" or "No", the logical deduction based on the provided numbers leads to "No". If the question implies a causal interpretation where smoking causes effort which causes admission, the math confirms the positive association. If the question is a trick regarding the base rate fallacy, the calculation above explicitly refutes the idea that smoking lowers admission chances.)*} \\
\\ 
\texttt{**Correction on the initial constraint:**} \\
\texttt{Upon re{-}evaluating the strict instruction "Start your answer with 'Yes' or 'No'", I must provide the correct logical answer derived} \\
            }
        }
    }
    & & \\ \\

    \theutterance \stepcounter{utterance}  
    & & & \multicolumn{2}{p{0.3\linewidth}}{
        \cellcolor[rgb]{0.9,0.9,0.9}{
            \makecell[{{p{\linewidth}}}]{
                \texttt{\tiny{[GM$|$GM]}}
                \texttt{Yes} \\
            }
        }
    }
    & & \\ \\

    \theutterance \stepcounter{utterance}  
    & & & \multicolumn{2}{p{0.3\linewidth}}{
        \cellcolor[rgb]{0.9,0.9,0.9}{
            \makecell[{{p{\linewidth}}}]{
                \texttt{\tiny{[GM$|$GM]}}
                \texttt{no} \\
            }
        }
    }
    & & \\ \\

    \theutterance \stepcounter{utterance}  
    & & & \multicolumn{2}{p{0.3\linewidth}}{
        \cellcolor[rgb]{0.9,0.9,0.9}{
            \makecell[{{p{\linewidth}}}]{
                \texttt{\tiny{[GM$|$GM]}}
                \texttt{game\_result = LOSE} \\
            }
        }
    }
    & & \\ \\

    \theutterance \stepcounter{utterance}  
    & & \multicolumn{4}{p{0.6\linewidth}}{
        \cellcolor[rgb]{0.95,0.95,0.95}{
            \makecell[{{p{\linewidth}}}]{
                \texttt{\tiny{[P1$\leftrightarrow$P]}}
                \texttt{Post{-}episode Reflexion (not sent to GM):} \\
\texttt{New plan: I failed because I calculated the conditional probabilities correctly but ignored the strict formatting constraint to start the answer with "Yes" or "No" before providing the reasoning. Although the logical answer was "No", I initially wrote "Yes" and then tried to correct myself mid{-}stream, which likely violated the "start with" rule or the expectation of a definitive initial stance. Next time, I will first perform the calculation to determine the correct boolean answer (Yes/No), and then immediately output that single word as the very first characters of my response, followed by the reasoning, ensuring no contradictory opening is made.} \\
            }
        }
    }
    & \\ \\

\end{supertabular}
}

\end{document}
